% vim: textwidth=78 ai spell spelllang=en
\documentclass{acm_proc_article-sp}
\usepackage[utf8x]{inputenc}
\usepackage[final]{hyperref}
\usepackage{color}
\usepackage{xspace}
\usepackage{graphicx}
\usepackage{listings}

\lstset{tabsize=2,showstringspaces=false,breaklines=true,breakatwhitespace=true,
	escapeinside={(*@}{@*)},basicstyle=\footnotesize}

\newcommand\Color[2]{\textcolor{#1}{#2}}
\newcommand{\nb}[1]{\marginpar{\scriptsize #1}}
\newcommand\todo[1]{\nb{\Color{red}{#1}}}

\newcommand\st{\textsc{Stanse}\xspace}

\begin{document}

\title{Eliminating Repetitive Errors Found in Consecutive Source Releases}

\numberofauthors{1}

\author{Ji\v{r}\'\i~Slab\'y\\
	\affaddr{Faculty of Informatics}\\
	\affaddr{Botanick\'{a} 68a}\\
	\affaddr{Brno, Czech Republic}\\
	\email{slaby@fi.muni.cz}
}

\maketitle

\begin{abstract}
With a use of contemporary static analysis tools, many errors are found in
source codes. We introduce a mechanism to prevent recurring of perpetual
errors found by the tools in different source releases. The mechanism aims
especially at those errors which were not fixed yet or are false-positives
which will likely be present with its source code forever. We discuss what
algorithms can be used to determine such errors and demonstrate them on output
from \st run on the Linux kernel.
\end{abstract}

\category{D.2.4}{Software}{Software Engineering}[Software/Program Verification]

\terms{Algorithms}

\keywords{Static analysis, error pruning, consecutive version errors,
iterative analysis}

\section{Introduction}
Recently, many static analysis tools emerged with a potential of finding many
long-standing errors in the source code. Some of the tools are open-source and
may be used for arbitrary projects~\cite{klee,java_bugfinding,sparse}, some
are commercial, but still check some picked open sources~\cite{coverity} or
are available for free use on developer's code~\cite{pex} and finally, some
are fully commercial~\cite{klocwork}.

These tools are based on a diverse set of algorithms, but in common, they find
errors in code, false-positives including; again with diverse ratio. When some
project which consists of million lines of code is checked by them, the error
list grows, because the project has more potential places for errors. And for
the same reason, more false-positives may be reported too.

When these projects are checked only once or even periodically but after a
long time, so that the sources changed much in the meantime and the reported
errors are different between the checking runs as well, there is a little
what can be done about the excessive error count. All of them need be
inspected (and fixed) manually.

On the other hand, when the project is checked on a regular basis, say daily
or weekly, the code is not changed that much. Then, the reported set of errors
almost does not vary and it becomes a burden for developers who would
have to check the very same subset of reported errors every day/week. And
given false-positives are not about to be fixed in the code at all, the chance
to avoid them is low except removing the files listed from the next to-check
list.

In the cases of repetitive analyses, when an analyzer is run on a project with
release cycle short enough, it would be convenient to recognize (a) the errors
which were reported in an earlier version and were not fixed
yet\footnote{There are many reasons for this including slow responses from
maintainers to error reports. An example of this is at
\url{http://lkml.org/lkml/2009/9/23/221} which had to be sent three times and
was merged even after 3 months.} and (b) false-positives because they are
persistent in time and nobody wishes to see them again and again in every
upcoming released version checks.

Although it may seem that there are almost no requirements on the project when
we want to design a technique dealing with (a) and (b), the assumption is not
strictly correct. Not only the fast development resulting in quick releases of
new versions is a precondition, but also rigid structure of files so that they
do not fluctuate over the source tree much. It was shown in~\cite{tracking},
used on Java code, that even a code fluctuation is not a blocker for this
approach and it can be handled without problems.

One of the projects satisfying the requirements above is the Linux kernel for
sure. Its release cycle between two major versions (e.g.
2.6.33$\rightarrow$2.6.34) is relatively short, oscillating between 11-13
weeks~\cite{linuxdev}. The changes coming into the tree are not accepted
unless appropriately reasoned, reviewed and merged by subtree maintainers and
finally accepted by L. Torvalds. Also, developers are mostly not allowed to
shuffle with sources arbitrarily over the tree, every class of drivers has its
own, predetermined, directory. Above that, changes without strong enough
reasoning are rejected.

We already mentioned that there are many static analysis tools out there. To
work with a fully operational tool capable to be applied to the Linux kernel,
\st~\cite{stanse}, a tool developed at the Faculty of Informatics, is used
specifically in this paper. It contains a base with C parser, error reporting
facility, user interface, points-to analysis module and application
programming interface (API) for checkers. A checker is a module which gets the
parser output and reports errors back to the base which in turn presents them
to a user. Currently, there are four checkers implemented, in this paper we
are interested exclusively in the \textsc{AutomatonChecker} which is able to
find errors in behaviour described by finite automata (such as in
Figure~\ref{fig:automaton}).

\begin{figure}[tb]
\begin{center}
\input{automaton.pspdftex}
\end{center}
\caption{Automaton for locking}
\label{fig:automaton}
\end{figure}

The paper is organized as follows. First, in Section~\ref{sec:motiv} we will
see why the technique is important for us and what benefits it brings, then
Section~\ref{sec:method} introduces the technique itself with possible
extensions and finally, we describe the implementation of the technique and
its evaluation in Section~\ref{sec:eval} with a use of
\st's output collected from the Linux kernel.

\begin{figure*}[tb]
\begin{center}
\begin{tabular}{|c|c|} \hline
\textbf{Older Version} & \textbf{Newer Version} \\ \hline \hline
\begin{lstlisting}[language=C]
static DEFINE_SPINLOCK(my_lock1);
static DEFINE_SPINLOCK(my_lock2);
static void fun(int in) {
  spin_lock(&my_lock1);
  if (in)
    return;
  spin_unlock(&my_lock1);
}
.
.
\end{lstlisting} &
\begin{lstlisting}[language=C]
static DEFINE_SPINLOCK(my_lock1);
static DEFINE_SPINLOCK(my_lock2);
static void fun(int a) {
  spin_lock(&my_lock1);
  if (a) {
    spin_unlock(&my_lock2);
    return;
  }
  spin_unlock(&my_lock1);
}
\end{lstlisting} \\ \hline
\end{tabular}
\end{center}
\caption{Illustration of false-negative introduced by the \emph{Simple
Solution}}
\label{fig:fp_simple}
\end{figure*}

\section{Motivating Example}
\label{sec:motiv}
Before we dive into the details of the algorithms in the
Section~\ref{sec:method}, consider the following piece of code extracted
(and also slightly manually macro-expanded for clarity) from
\texttt{lib/locking-selftest.c} obtained from the Linux kernel:
\begin{center}
\begin{lstlisting}[language=C]
static DEFINE_SPINLOCK(lock_A);
...
static void double_unlock_spin(void) {
  spin_lock(&lock_A);
  spin_unlock(&lock_A);
  spin_unlock(&lock_A);
}
\end{lstlisting}
\end{center}

The \texttt{double\_unlock\_spin} function intentionally breaks the proper
locking policy expressed by the automaton in Figure~\ref{fig:automaton}. The
\texttt{lock\_A} is locked and then unlocked twice, thus leaving the lock in
an inconsistent, erroneous state (\emph{Error} state). However, in this case,
it is a sole purpose of the function. It is intended for exercising the
dynamic mechanism for automatic locking violations detection present in the
kernel\footnote{So-called \textsc{LockDep}, more information is available in
the article at \url{http://lwn.net/Articles/185605/}.}. Thus it is impossible
to change or "fix" the code in any way to be correct from the point of view of
the automaton in Figure~\ref{fig:automaton}. Hence it is an error (a kind of
false positive) which will be with this part of sources forever.

We investigated how much this case occurs while bugs are being found.  We
started with a complete review of errors reported in the Linux 2.6.27 and
sorted out which of those are false-positives. They were marked accordingly
and true errors reported to corresponding code maintainers (some with
appropriate patches). Then, the analysis started being performed more often,
even on release candidate kernels, approximately in a weekly period. When
inspecting these versions, we found out that many of the errors repeats.
We investigated that they are of the similar kind like the example above.
I.e. they cannot be fixed (false positive) or were already reported but not
fixed yet for some reason (too slow communication, lack of response, \ldots).

To tackle down this problem, we present an algorithm to match errors from
consecutive code releases to avoid repetition of errors given to developers in
error lists. Thus lowering the burden exposed to them when walking through the
error list generated from later versions. Clearly, every reported error from
each iteration has to be fixed, otherwise it is lost once for ever. To the
best of our knowledge there is no white paper describing a similar technique.

\section{Methodology}
\label{sec:method}
Before we start with the methodology, we first have to know what the actual
output of the tools look like, especially what the \st's output is. Even
though the error formatting practices differ tool by tool, in common they
contain the source file name where the error is located with exact line too.
The report also carries an information about what kind of error is that and
possibly what module (checker) generated that. Eventually some tools report
another fields like line column or relevance, but this is irrelevant for
further considerations, because in principle, the extension can be done
easily.

We consider a list of reports $RL$ to be a set of entries (errors) $E$. The
error is represented as a tuple for reasons stated later and looks like
in (\ref{eq:edef}). \textit{file} and \textit{line} both correspond to the
location of the error found by \textit{checker} and described in
\textit{error\_description}.
\begin{equation} \label{eq:edef}
E = (file, line, checker, error\_description) \in RL
\end{equation}

Now, we can have a starting point two lists of errors (sets of tuples) $RL_1$
and $RL_2$ generated from two distinct versions of the same project. $RL_1$
corresponds to an older version, $RL_2$ to a newer one respectively.

Having the lists, we want to eliminate such tuples $E \in RL_2$ which
correspond to the same error in $RL_1$. The "same error" term is not precisely
defined here. The following subsections will define it differently, based on
actual needs. Finally, we note the result of elimination as $RL^{filtered}_2$.

\subsection{Simple Solution}
\label{sec:simple}
The most straightforward approach of defining sameness is just to compare each
tuple from $RL_2$ to entries in $RL_1$ and obtaining the result
(\ref{eq:simple}).
\begin{equation} \label{eq:simple}
RL^{filtered}_2=RL_2\setminus RL_1
\end{equation}
%\{x \mid x\in RL_2 \wedge x \notin RL_1\}

Obviously there are several problems tightly connected to this approach.
First, not all the tuple elements are preserved over source revisions. Namely,
an example is \textit{line} component changing extensively version by version.
To make this technique defunct, it is actually enough that somebody adds or
removes a single (even empty) line to/from the file which the error is
reported in\footnote{Technically, if the analyzer works on the preprocessed
code, it is enough that somebody changes any included file, because the
preprocessed output is shifted by the change as a whole.}.

Second, this solution can introduce false-negatives. Consider a code in
Figure~\ref{fig:fp_simple} on the left side where there is a bug in locking.
The function \texttt{fun} omits to unlock \texttt{my\_lock1} depending on the
input parameter \texttt{in}. This error is reported to the maintainer of the
code and they try to fix it. Unfortunately, they change the code wrong (the
right side of the Figure) and add a new error caused by unlocking a wrong
lock.

Normally, as the error differs in no way from any other one, this would not
matter -- the analyzer can indeed find and report it in every check performed.
However with the eliminating technique enforced, a problem emerges. It stems
out of the error being on the same line like the one reported in the previous
release. So in the end for both of the errors the entry will look the same in
report lists from both runs:
\begin{displaymath}
E=(file.c, 6, AutomatonChecker, Locking\_imbalance)
\end{displaymath}
Inconveniently, the outcome of this is that the second error is hidden from
programmer when the \emph{Simple Solution} is used because $E \in RL_1$ and
hence, from the equation~(\ref{eq:simple}), the error will not appear in
$RL^{filtered}_2$ anymore.

\begin{table*}[tb]
\begin{center}
\begin{tabular}{|l|r||r|r||r|r||r|r|} \hline
          & \textbf{none} & \textbf{simple} & \textbf{FR} (\%) &
	  \textbf{brutal} & \textbf{FR} (\%) & \textbf{advanced} &
	  \textbf{FR} (\%) \\ \hline
Atomicity        & 23  & 11  & 52.2 & 0  & 100.0 & 11  & 52.2 \\ \hline
Locking          & 179 & 64  & 64.3 & 5  & 97.3  & 64  & 64.3 \\ \hline
Memory           & 217 & 55  & 74.7 & 11 & 95.0  & 55  & 74.7 \\ \hline
Pairing          & 81  & 11  & 86.5 & 0  & 100.0 & 11  & 86.5 \\ \hline
\textbf{Summary} & 500 & 141 & 71.8 & 16 & 96.8  & 141 & 71.8 \\ \hline
\end{tabular}
\end{center}
\caption{Error count using none, simple, brute-force and advanced solutions
respectively. \emph{FR} stands for \emph{Filter Ratio} and is explicitly
stated for all methods.}
\label{tab:res}
\end{table*}

\subsection{Advanced Solution}
\label{sec:adv}
The issue noted above is not the only introduced by the \emph{Simple
Solution}. Before we solve that particular one, we will look into one more
issue caused by \textit{file} component and we will also introduce a
supporting foundation for dealing with the problems.

Both \textit{line} and \textit{file} changes are caused by code shuffling, the
former because of code reordering (e.g. to define a function or variable prior
their use without forward declarations), the latter because of file
moves/renames (subsystem restructuring etc.). Obviously, the problem still
remains to be that we need to match \textit{line} and \textit{file} components
of an error tuple $E \in RL_2$ to tuples in $RL_1$.

We can achieve \textit{file} to match by modern \emph{Software Configuration
Management} (SCM) that is able to track renames. A perfect example of such
system is \texttt{git} \cite{git}. When the two releases are given as tags,
the files changed are output and above that, renamed files marked with
"\texttt{=>}" when special option passed to the SCM (\texttt{-M} in the
\texttt{git} case). For example when \texttt{old.c} is renamed to
\texttt{new.c} between \texttt{oldrev} and \texttt{newrev} revisions, the
changes output then looks like the following:
\begin{lstlisting}[language=sh]
$ git diff --stat -M oldrev..newrev
1       0       old.c => new.c
\end{lstlisting}

Having this output, for our purposes, we silently ignore the numbers at the
beginning and from the rest, we extract two sets $S_{old}$ and $S_{new}$
containing the old and new filenames respectively. We also have a renaming
function $r_f$ which accepts a tuple as parameter such that
$r_f(("old.c",l,c,e))="new.c"$ for each found change ($"old.c" \in S_{old}$,
$"new.c" \in S_{new}$). In contrary, if the filename is unchanged, i.e.  $x
\notin S_{old}$, the function is defined as $r_f(x)=first(x)$ (return the
first component of $x$).

Then we extend the function $r_f$ for tuples component-wise while using the
same approach on other components of the tuple. For the time being, they are
defined as $r_l=second$, $r_c=third$ and $r_e=fourth$ (just pick up the
correct component):
\begin{displaymath}
new_1(x) = (r_f(x),r_l(x),r_c(x),r_e(x))
\end{displaymath}
Now, we define $RL_1^{new}$ in~(\ref{eq:rl_file}) and replace $RL_1$ by that
in equation~(\ref{eq:simple}). Thus, written informally, we approach the error
list from the older run $RL_1$ to the newer one $RL_2$ so that they can be
compared easier.
\begin{equation} \label{eq:rl_file}
RL_1^{new} = \{new_1(x) \mid x \in RL_1\}
\end{equation}

The \textit{line} case described in the previous section may be worked out
similarly, but even with more machinery, because $RL_2$ needs to be changed as
well. So similarly to $RL_1^{new}$, we define $RL_2^{new} = \{new_2(x) \mid x
\in RL_1\}$ where $new_2(x) = (s_f(x),s_l(x),s_c(x),s_e(x))$. In this case,
$s_f$, $r_c=s_c$, $r_e=s_e$ are defined as pick-up functions like above. $r_f$
remains the same from the previous definitions to have file renaming
detection. The remaining $r_l$ and $s_l$ are defined to return 2-tuples as:
\begin{displaymath}
\begin{matrix}
r_l((f,l,c,e))=(l, fetchline(oldrev,f,l)) \\
s_l((f,l,c,e))=(l, fetchline(newrev,f,l))
\end{matrix}
\end{displaymath}

$fetchline$ there loads the line $l$ from the file $f$ of the corresponding
revision. Then, $RL_2^{new}$ is used instead of $RL_2$ and the filtered set is
defined as:
\begin{displaymath}
RL_2^{filtered} = \{x \mid x \in RL_2^{new} \wedge x \notin RL_1^{new}\} \\
\end{displaymath}
with typical entries like:
\begin{displaymath}
E=(file,(line,line\_contents),checker,error\_description)
\end{displaymath}

At this point, it remains to reason, how comparison of the line contents with
line number itself avoids the drawn problem. First, the errors are reported on
lines which are not empty, they rather contain some meaningful code. This
ensures, that there will be something to check against. Second, if (line, line
contents) tuple is the same, it is an error which was reported already and
should be in the process of fixing, i.e.  we do not include it in the filtered
set. Finally, when the code on the line is not the same, there is a new bug
(like in Figure~\ref{fig:fp_simple}). The \emph{Advanced Solution} is hence
supposed to leave us on the safe side, so that no real bugs are left
unattained after some period of time.

To some extent, all the functions described here are defined generically. For
example, we did not define $r_c$, $r_e$, $s_c$ and $s_e$ in our computations
in some thorough way at all, but they may be redefined arbitrarily depending
on concrete needs, including the case of different definitions of tuples.

\section{Evaluation}
\label{sec:eval}
We exercised our algorithm on the errors found in the Linux kernel by \st. The
two versions we considered are 2.6.35-rc3-next20100618 and
2.6.35-rc4-next20100706\footnote{-next is a special development branch with
all unstable features being merged every day.}, further referred only as -rc3+
and -rc4+ respectively. That means development versions where there are over
300 thousand lines changed in 2654 files between them in 18 days. We made no
changes in configuration of the kernel (besides newly added drivers, which
were ignored) and neither of our tool.

The errors reported by \st are represented in an XML form by default. But as
we wanted to have a convenient way for combining and querying the data, we
decided to use databases and some SQL language for accessing the data. The
reason is also to easily represent the tuples. This way, a tuple is simply a
row in a database table.

Although there already exist infrastructures for combining XML and SQL like
SQL/XML \cite{sqlxml}, we sticked to a simple database system, more concretely
to \textsc{SQLite} \cite{sqlite} which stores and references databases as
single files somewhere on a backing store.

To achieve the goal, we wrote a Perl script to transform the XML into an
\textsc{SQLite} database file. Despite each report list resides in a separate
database file now, it is not a problem because \textsc{SQLite} allows
combining even tables from different databases in its SQL. Thus we can easily
reference errors from two versions and store the result into the third as we
will see later.

Having separate databases for each list allowed us fast deployment in our web
frontend too, because it works on a version-per-file basis and needed no
extra changes with this approach.

\subsection{Using Simple Solution}
First, we will have a look at the \emph{Simple Solution} presented in
Section~\ref{sec:simple}. Overall, the results are summarized in
Table~\ref{tab:res}. The \emph{Atomicity} (e.g. forbidden sleeps inside atomic
context), \emph{Locking} (omitted unlocks, double locks, etc.),
\emph{Memory} (NULL dereferences, dangling pointers and so on) and
\emph{Pairing} (grabbed resource not given back to the system for example)
rows correspond to the errors found by the
\textsc{AutomatonChecker} with different configuration files \footnote{The errors found by \st (with their
description) may be seen at \url{http://stanse.fi.muni.cz/bugs.html}.}.

We also expressed the reduction in percents as \emph{Filter Ratio}.  It
designates how many errors were pruned and is computed the standard way as
$100-\frac{100\cdot |RL_2^{filtered}|}{|RL_2|}$. There are three columns
containing them, the first when -rc4+ with \emph{Simple Solution} is
considered, the last for \emph{Advanced Solution} being used. The middle one
is described later.

Unexpectedly, they demonstrate, that it does matter on which type of errors
the technique is used; the filter ratio oscillates much. We have no direct
explanation for that. After we walked through the errors, we sort out that no
pattern can be deduced from them. However we can only think that there is none
at best.

For this solution, the SQL query, sort and insert command executed in the
database which would contain the filtered set is as below. The two report
lists are attached as \emph{new} and \emph{old} and used in the query. The ||
token there stands for string catenation.
\begin{lstlisting}{language=SQL}
ATTACH DATABASE "35-rc3.db" AS old;
ATTACH DATABASE "35-rc4.db" AS new;
INSERT INTO errors (checker, importance, error, file, line)
  SELECT checker, importance, error, file, line FROM new.errors WHERE
  file || line || checker || error NOT IN
    (SELECT file || line || checker || error FROM old.errors);
\end{lstlisting}

The table also contains a column for brute-force approach which is even more
naive than the \emph{Simple Solution}. The query above is changed so that the
line number is completely removed. In this case we may lose some new bugs,
i.e.  have even more false-negatives than earlier, but it still seems to be
enough in cases where \st is run after each -next kernel release -- on a daily
basis and that is exactly the use case of this variant.

Either way, we can see that we eliminated 359 (71.8\,\%) and with brute-force
even 484 (96.8\,\%) errors which are out of developers' attention now.

\subsection{Using Advanced Solution}
We present the results from the same setup along with the previous ones in
Table~\ref{tab:res} too. For our testing sample, the numbers are exactly the
same. This solution's primary aim is to generate safe results, rather than
better results. So this means that in these two particular versions, the
\emph{Simple Solution} generated safe enough results which technically needs
not be always the case.

We have not seen such behaviour yet though. Besides code changes, it it also
because this threat is mitigated with our tool due to how wrong the
preprocessed output is taken. It is not taken as a whole, but includes, which
might change the line numbers more often are not a part of the checking
process, so they are not involved in the reported lines.

On the other hand, it by no means renders this solution useless. After we fix
up the parser in \st, and consider whole sources, the situation might change.
The same way as when instead of \st, some other tool is used.

Leaving the usefulness apart, the script used previously for converting XML to
database was augmented to fetch and store also the line contents from the file
into the database. The structure of the database table was altered a bit to
contain one more field.  So, instead of storing 2-tuple of line and its
contents in one column, they are stored in two separate ones.

The SQL query used to obtain the filtered results here is just an enhancement
of the former -- \textit{line\_contents} is added there to appropriate places.

\newpage

\section{Conclusion}
We presented an algorithm for pruning errors which are reported repeatedly
from consecutive versions of a software project by static analysis tools. We
introduced two variants of the algorithm that may be used and we reasoned what
it can lead to when only the simpler one is deployed. Generally, given two
different versions, the technique renders an arbitrary analysis into the
iterative version.

The method was implemented and then evaluated on the Linux kernel. We found
out that the need of walking through the errors is rapidly lowered with the
algorithm by almost three quarters and thus the person-hours needed otherwise
are reduced. Hence potentially the overall cost of the product being checked.

\section*{Acknowledgements}
The author is supported by the research centre Institute for Theoretical
Computer Science (ITI), project No. 1M0545. I thank my colleagues for their
valuable ideas.

\bibliographystyle{plain}
\bibliography{inter-errors}

\end{document}
